\centerline{\bf \Huge pqr-метод}

В этом файле мы постараемся объяснить суть популярного метода решения симметрических неравенств. В первую очередь мы начнём с некоторого утверждения.

\textbf{Утверждение.} 
По тройке действительных чисел $a, b, c$ можно однозначно определить тройку чисел

$$
p=a+b+c, \quad q=a b+b c+c a, \quad r=a b c
$$

Аналогично по тройке чисел $p,q,r$ можно восстановить тройку $a,b,c$. 

Первая часть утверждения очевидна. Чтобы восстановить по тройке $p, q, r$ числа $a, b, c$, необходимо решить уравнение $x^3-p x^2+q x-r=0$. \textbf{Именно с этим многочленом мы и будем работать}. Однако у этого уравнения не всегда все корни действительные (например, уравнение $x^3 -1=0$ имеет два комплексных корня). В случае квадратного уравнения у нас был чёткий критерий того, что оба корня действительные --- дискриминант должен быть неотрицательным. Оказывается, что это утверждение обобщается на случай многочлена произвольной степени.

\textit{Определение.}
Для многочлена $P(x)=a_0+a_1 x+\cdots+a_n x^n, a_n \neq 0$, его \textbf{дискриминант} есть выражение:

$$
\mathcal{D}(P)=a_n^{2 n-2} \prod_{i < j}\left(x_i-x_j\right)^2
$$

где $x_1, x_2, \ldots, x_n$ --- все корни многочлена (с учётом кратности).

\textbf{Критерий вещественности.}
Многочлена $P(x)$ имеет вещественные корни тогда и только тогда, когда $D(P)\geq 0$

В нашем же случае дискриминантом многочлена $x^3-p x^2+q x-r=0$ будет равен

$$
\mathcal{D} = (a-b)^2(b-c)^2(c-a)^2
$$

Спустя пару часов страданий можно выяснить, что $D =-4 p^3 r+p^2 q^2+18 p q r-4 q^3-27 r^2=T(p, q, r)$

\newpage

\hspace{3mm}

В задачах на неравенства переменные обычно неотрицательные. Поэтому от переменных $a,b,c$ хочется попросить не только вещественности, но и не отрицательности. 

\textit{Определение.}
Назовём тройку $(p, q, r)$ \textbf{допустимой}, если ей соответствует тройка действительных неотрицательных чисел $a, b, c$.

\textbf{Критерий неотрицательности/допустимости.}
Тройка $(p, q, r)$ допустима тогда и только тогда, когда

$$
p \geqslant 0, q \geqslant 0, r \geqslant 0, T(p, q, r) \geqslant 0 .
$$

Если за неравенствами $p \geqslant 0, q \geqslant 0, r \geqslant 0$ проследить ещё легко, но вот за неравенством $T(p,q,r)\geq0$ проследить довольно трудоемко, так как мы не совсем понимаем, как себя ведёт функция от трёх переменных, график которой, вообще говоря четырёхмерный. Поэтому мы будем фиксировать какие-то две переменные и следить только за третьей. Следующие три леммы нам расскажут как правильно это делать.


\textbf{Лемма об $r$.}
Пусть для заданных $p=p_0$ и $q=q_0$ существует хотя бы одно $r$ такое, что тройка $(p_0, q_0, r)$ допустима. Докажите, что для минимального $r$ такого, что тройка $(p_0, q_0, r)$ допустима, в соответствующей тройке корней или есть два равных, или один из них равен 0 . А для максимального $r$ такого, что $(p_0, q_0, r)$ допустима, верно, что в соответствующей тройке корней есть два равных.

\textbf{Лемма о $q$.}
Пусть для заданных $p=p_0$ и $r=r_0$ существует хотя бы одно $q$ такое, что тройка $(p_0, q, r_0)$ допустима. Докажите, что для минимального и максимального $q$ такого, что $(p_0, q, r_0)$ допустима, в соответствующей тройке корней есть два равных.

\textbf{Лемма о $p$.}
Пусть для заданных $q=q_0$ и $r=r_0>0$ существует хотя бы одно $p$ такое, что тройка $(p, q_0, r_0)$ допустима. Докажите, что для минимального и максимального $p$ такого, что $(p, q_0, r_0)$ допустима, в соответствующей тройке корней есть два равных.

Выглядят эти леммы довольно устрашающе и не совсем понятно, как они нам могут помочь. Давайте разберёмся на примере леммы об $r$ и леммы о $q$.

\newpage

\hspace{3mm}

В лемме об $r$ мы зафиксировали какие-то значения $p$ и $q$ и двигаем значение $r$. Тогда многочлен $T(p,q,r) = -4 p^3 r+p^2 q^2+18 p q r-4 q^3-27 r^2$ теперь становится многочленом от одной переменной, причем второй степени. Учитывая, что старший коэффициент отрицательный парабола будет направлена вниз.

На рис.1 зелёным отображена та область, в которой дискриминант будет неотрицательным, то есть именно та область значений $r$, в которой мы можем работать.

\begin{figure}[ht]
  \centering
  \includegraphics[width=.6\linewidth]{Теория/pqr_pics/quad_1.png} 
  \caption{Левый корень $r=0$}
\end{figure}

Заметим, что она всегда будет находиться между корнями параболы, так как старший коэффициент отрицательный. Стало быть минимальное и максимально допустимое значение $r$ это ни что иное как корни уравнения $T(p_0,q_0,r)=0$. Но расположение корней может быть совсем разным, например, на рис. 1 левый корень равен нулю, а правый положительный. Так же может быть ситуация, когда левый корень отрицательный (см. рис. 2).

\newpage

\hspace{3mm}

В такой ситуации мы должны помнить, что переменные $p,q,r$ неотрицательные, поэтому всё равно левое допустимое значение $r$ будет равно нулю, а правое так же положительное. Ну и последняя ситуация, когда оба корня положительные (оба отрицательные нас не устраивает, так как $r$ должна принимать хотя бы одно неотрицательное значение).

\begin{figure}[ht]
  \centering
  \includegraphics[width=.5\linewidth]{Теория/pqr_pics/quad_2.png} 
  \caption{Левый корень отрицательный}
\end{figure}

Какой из этого всего можно сделать вывод? Крайние допустимые значения $r$ это корни уравнения $T(p_0,q_0,r) =0$ за исключением того случая, когда левый корень отрицательный. В этом случае минимальное допустимое значение $0$, а максимальное это какой-то положительный корень выше упомянутого уравнения. Но что нам даёт наличие корня уравнения $T(p_0,q_0,r) =0$? Давайте вспомним, что $T(p,q,r) = (a-b)^2(b-c)^2(c-a^2)$ и обнуляется это выражение тогда и только тогда, когда какие-то двое из тройки чисел $a,b,c$ равны. А если же $r=0$, то $abc=0$, а значит какое-то из чисел равняется $0$. Теорема доказана.


\newpage

\hspace{3mm}


Перейдём к теореме о $q$. Тут ситуация похожая, но теперь при фиксированных $p$ и $r$ многочлен $T(p,q,r) = -4 p^3 r+p^2 q^2+18 p q r-4 q^3-27 r^2$ будет третьей степени, у которого старший коэффициент отрицательный. У такого многочлена либо все три корня отрицательные, либо один отрицательный и два положительных. Нас устраивает только лишь последний вариант (см. рис. 3)

\begin{figure}[ht]
  \centering
  \includegraphics[width=.6\linewidth]{Теория/pqr_pics/cube.png} 
  \caption{Два корня положительны, один отрицателен}
\end{figure}

Стало быть минимальное и максимальное допустимое $q$ это два строго положительных (так как свободный член отрицательный) корня. А мы уже знаем, что в корнях какие-то два числа из тройки $a,b,c,$ будут равны. Это заканчивает доказательство теоремы. Теорема о $p$ доказывается аналогично теореме о $q$.

\newpage

\hspace{3mm}


В чём же \textbf{идея метода}? Пусть дано симметрическое неравенство от трёх переменных. Сведём его к симметрическому неравенству от $p, q, r$. Зафиксируем два из чисел $p, q, r$ и минимизируем/максимизируем оставшееся число. Тогда исходное неравенство сведётся к неравенству от двух переменных.

\textbf{Пример 1.}
Для неотрицательных чисел $a, b, c$ докажем неравенство

$$
a^2+b^2+c^2 \geqslant a b+b c+c a .
$$


Перепишем его в виде $p^2-2 q \geqslant q \Leftrightarrow p^2 \geqslant 3 q$. Зафиксируем, к примеру, значения $p=p_0$ и $r=r_0$. Теперь надо доказать неравенство $q \leqslant \frac{p_0^2}{3}$, то есть показать, что $q$ ограничено сверху некоторой константой. Заметим, что если мы докажем неравенство для максимального $q$, то оно будет верно для всех $q$. По лемме о $q$ максимальное значение $q$ достигается, когда какие-нибудь из переменных $a, b, c$ равны. Без ограничения общности, $a=b=x, c=z$. Поскольку неравенство $q \leqslant \frac{p_0^2}{3}$ равносильно исходному, то после подстановки в исходное неравенство получим

$$
2 x^2+z^2 \geqslant x^2+2 x z \Leftrightarrow(x-z)^2 \geqslant 0,
$$


что верно. Следовательно, неравенство верно для всех положительных чисел $a, b, c$.

\textbf{Пример 2.}
В этом примере мы разберем один полезный трюк связанный с переменной $r$.

Неотрицательные числа $a, b, c$ таковы, что никакие два из них не равны нулю. Докажите, что

$$
\frac{1}{a^2+b^2}+\frac{1}{b^2+c^2}+\frac{1}{c^2+a^2} \geqslant \frac{10}{(a+b+c)^2}
$$


Решение. Исходное неравенство переписывается следующим образом:

$$
\frac{p^4-4 p^2 q+5 q^2-2 p r}{q^2 p^2-2 q^3-2 p^3 r+4 p q r-r^2} \geqslant \frac{10}{p^2}
$$

$$
\left(p^4-4 p^2 q+5 q^2-2 p r\right)-\frac{10}{p^2}\left(q^2 p^2-2 q^3-2 p^3 r+4 p q r-r^2\right)\geq 0
$$

$$
\frac{10}{p^2} r^2+\left(18 p-\frac{40 q}{p}\right) r+\left(p^4-4 p^2 q-5 q^2+\frac{20 q^3}{p^2}\right) \geqslant 0
$$


Коэффициент при $r^2$ положительный следовательно, минимальное значение функция принимает при минимальном $r$. По теореме об $r$ мы знаем, что минимальное значение либо $0$, либо какой-то корень уравнения $T(p,q,r)=0$. Рассмотрим два случая:

1) Если минимальное значение положительное, то $a=b$. \textbf{Теперь собственно фишечка}. Обозначив $t=\frac{a}{c}>0,a = b = tc$), исходное неравенство внезапно теперь будет зависеть от одной переменной:

$$
20 t^3-11 t^2+4 t+1 \geqslant 0
$$

Но почему так? Давайте подставим вместо чисел $a,b,c$ числа $ta, ta, tc$ в исходное неравенство и посмотрим, что получится:

$$
\frac{1}{(t c)^2+(t c)^2}+\frac{1}{(t c)^2+c^2}+\frac{1}{c^2+(t c)^2} = \frac{10}{(t c+t c+c)^2}
$$

сокращаем обе части на $\dfrac{1}{c^2}$ и получаем выражение зависящее только от $t$:

$$
\frac{1}{2 t^2}+\frac{2}{1+t^2} \geq \frac{10}{(2 t+1)^2}
$$

Ну а дальше домножаем на общий знаменатель, приводим подобные слагаемые  и получаем то самое неравенство $20t^3-11t^2+4 t+1 \geqslant 0$, которое выполнено для всех положительных $t$ и побеждаем. Волшебное сокращение на $\dfrac{1}{c^2}$ связано с тем, что исходное неравенство было \textbf{однородным} (вспомните определение однородного неравенства), а поэтому при замене $a=b=tc$ мы смогли вынести одну и ту же степень $c$ и сократить.


\newpage

\hspace{3mm}

2) Теперь, если минимально допустимое значение $r$ равно нулю, то $a=0$. Опять делаем фишечную замену $t=\frac{b}{c}>0, a = b = tc$ и перепишем исходное неравенство в виде:

$$
(t-1)^2\left(t^4+4 t^3+t^2+4 t+1\right) \geqslant 0,
$$

что очевидно верно при положительных $t$